% ===================================================================
%  TABULKA č. 10 — Moře. — Přístav.
%  Traduction 2026 d'apres texto/io, controlee sur texto/fr, texto/ru
%  et texto/uk. Consigne : voir texto/cs/10-tabulka-01.tex.
%
%  Le francais decale sa numerotation d'un cran a partir de l'alinea
%  8. C'est une coquille de l'imprimeur francais ; le tcheque suit
%  l'ido.
%
%  LE TABLEAU DE LA MER EST LE PLUS EMPRUNTE DE TOUTES LES COLONNES,
%  ET LE TCHEQUE Y AJOUTE UN CAS QUE LE VOLUME N'AVAIT PAS ENCORE
%  RENCONTRE : c'est la premiere langue de la serie qui n'a pas de
%  mer. La Boheme est enclavee depuis toujours, et sa marine est un
%  vocabulaire de LIVRE, non de port. « Obrněnec », « torpédoborec »,
%  « křižník », « fregata », « šroub », « kýl », « kormidlo »,
%  « maják », « bóje », « přístav » : tout cela existe et se dit,
%  mais rien n'y est venu par un quai.
%
%  ET CELA SE VOIT A UN DETAIL QUI NE SE VOIT NULLE PART AILLEURS.
%  La ou l'irlandais et le galicien opposaient la marine
%  internationale aux mots de la COTE — le brisant, la falaise,
%  l'ecueil — le tcheque n'a pas de lexique de cote a opposer. Il
%  emploie « útes » pour la falaise ET pour le recif, la ou le
%  galicien avait « cantil » et « escollo », et « úskalí » — de
%  « skála », le rocher de montagne — la ou l'irlandais avait
%  « sceir ». La regle du lieu d'apprentissage donne ici son cas
%  limite : la ou la chose n'est jamais venue, il n'y a ni emprunt ni
%  mot propre, il y a une TRADUCTION.
% ===================================================================

%%K t10-apar-1 apar
\VUsaut{4.0mm}
\VUtitre{15.0pt}{60}{TABULKA č. 10}
\VUsaut{3.0mm}
\VUpk{11.4pt}{40}{Moře. --- Přístav.}
\VUsaut{3.0mm}

%%K t10-01-1 p
1. --- V létě, zatímco jedni hledají klid a chlad na horách, jiní dávají přednost cestě k
moři. Tak jsme to udělali loni, protože máme velmi rádi \VUgras{oceán} \textsuperscript{(1)}.

%%K t10-02-1 p
2. --- Co se mne týče, mám velké potěšení z pozorování té nesmírné vodní plochy, která se
rozprostírá až k \VUgras{obzoru} \textsuperscript{(2)}, a z pohledu na obrovské
\VUgras{parníky} \textsuperscript{(3)}, jež odplouvají na dlouhou \VUgras{plavbu} s
cestujícími mířícími do Ameriky.

%%K t10-03-1 p
3. --- Proto můj otec, který zná mé záliby, vybírá vždy k strávení prázdnin velmi klidnou pláž
ležící blízko jednoho z našich velkých \VUgras{válečných přístavů.}

%%K t10-03-2 p
Za našeho posledního pobytu \VUgras{eskadra} přijela zakotvit za obrovskou
\VUgras{hrází} \textsuperscript{(4)}, která uzavírá a chrání \VUgras{vojenský přístav}
\textsuperscript{(5)}, a podařilo se mi podle přání obdivovat různé \VUgras{válečné lodi}:
malou \VUgras{ponorku} \textsuperscript{(10)}, která se noří a mizí pod vodou, i obrovský
\VUgras{obrněnec} \textsuperscript{(9)}, který se doposud bál téměř jen útoků
\VUgras{torpédovky} \textsuperscript{(8)}.

%%K t10-tit-2 sub
\VUsaut{3.0mm}
\VUpk{10.2pt}{40}{Malé lodi.}
\VUsaut{3.0mm}

%%K t10-04-1 p
4. --- Ta už uměla velmi \VUgras{obratně} najít slabé místo tlustého kovového pancíře, aby tam
zabodla nebezpečné \VUgras{torpédo}, jehož výbuch způsobí zkázu lodi; ale přece byla méně
obávaná než ponorka, protože torpédoborce ji snadno pronásledují.

%%K t10-05-1 p
5. --- Mohl jsem také vidět rychlé obrněné \VUgras{křižníky} \textsuperscript{(6)}, téměř tak
nezranitelné jako pravý obrněnec, několik \VUgras{dopravních lodí} \textsuperscript{(12)}, tři
nebo čtyři \VUgras{dělové čluny} \textsuperscript{(7)} a konečně \VUgras{fregatu}
\textsuperscript{(11)}.
\VUgras{Podívaná na to loďstvo, když manévruje, aby zvedlo kotvy, je ze všeho nejzajímavější}.

%%K t10-06-1 p
6. --- Jaký rozdíl ve stavbě mezi těmi velkými ocelovými hmotami a elegantními
\VUgras{trojstěžníky} nebo půvabnými \VUgras{plachetnicemi} \textsuperscript{(13)},
dosud užívanými v obchodním loďstvu, které jsem viděl vplouvat do přístavu se všemi plachtami,
když jsem se procházel po \VUgras{molu} \textsuperscript{(14)}. Ovšem, tyto ztrácely mnoho ze
svých výhod, \VUgras{když} byl \VUgras{vítr} nepříznivý. Musely stáhnout všechny plachty, aby
znovu vpluly do bazénu, a \VUgras{vlečný parník} \textsuperscript{(16)} byl nutný, aby je
vyzvedl a přivedl, jako prosté \VUgras{nákladní čluny} \textsuperscript{(17)}, k nábřeží.

%%K t10-tit-3 sub
\VUsaut{3.0mm}
\VUpk{10.2pt}{40}{Obchodní přístav.}
\VUsaut{3.0mm}

%%K t10-07-1 p
7. --- Co je zajímavé na přístavu, o němž mluvím, je to, že neobsahuje jen vojenský přístav,
uměle vytvořený obrovskými pracemi, ale také podivuhodný \VUgras{obchodní přístav} položený v
\VUgras{ústí} velké řeky.

%%K t10-08-1 p
8. --- Jazyk země, který odděluje oba bazény, je opevněn a hájen řadou \VUgras{baterií} a
\VUgras{bastionů}, které vystupují do moře. Je třeba zvláštního povolení k procházce po
\VUgras{hradbách} \textsuperscript{(15)}. Snadno jsem je získal prostřednictvím svého strýce,
který je inženýrem \VUgras{loděnic} \textsuperscript{(18)}, a šel jsem se podívat na lodi,
které se stavějí pod prostornými kůlnami.

%%K t10-09-1 p
9. --- Byl jsem dokonce přítomen spuštění obrněnce na vodu a vzpomínám si na dojemný okamžik,
který pocítil celý zástup, když se obrovská hmota, ponechána sama sobě, dala do klouzání po
své skluznici a připlula na vodu řeky.

%%K t10-10-1 p
10. --- Aby se člověk dostal do loděnic, přejde \VUgras{cvičiště} \textsuperscript{(19)}, kam
vojáci posádky přicházejí denně cvičit, a přejde po železné \VUgras{lávce}
\textsuperscript{(20)}, která spojuje přehlídkovou esplanádu s
\VUgras{arzenálem} \textsuperscript{(21)}.

%%K t10-tit-4 sub
\VUsaut{3.0mm}
\VUpk{10.2pt}{40}{Arzenál a doky.}
\VUsaut{3.0mm}

%%K t10-11-1 p
11. --- Se svým strýcem jsem navštívil tu velkou budovu plnou \VUgras{děl}
\textsuperscript{(22)}, \VUgras{koulí} \textsuperscript{(23)} a \VUgras{granátů}. Viděl jsem
také \VUgras{hutní továrnu} \textsuperscript{(24)}, v níž se vyrábějí
\VUgras{kotle} \textsuperscript{(25)} parníků.

%%K t10-12-1 p
12. --- Velmi blízko jsou \VUgras{doky} \textsuperscript{(26)} --- opravárenské doky --- a
velmi pozoruhodné \VUgras{plovoucí doky} \textsuperscript{(27)}, v nichž jsem mohl vidět zcela
zblízka, na opravované lodi, \VUgras{lodní šroub} \textsuperscript{(28)}, \VUgras{kýl}
\textsuperscript{(29)} a \VUgras{kormidlo} \textsuperscript{(30)} člunu.

%%K t10-13-1 p
13. --- Obchodní přístav stojí za prostudování zrovna tak jako vojenský přístav, a strávil
jsem mnoho hodin zevlováním na nábřežích, kde jsou přivázány \VUgras{kolesové parníky}
\textsuperscript{(31)}, které plují proti proudu řeky, a pozorováním, jak pracuje
\VUgras{stěžňový jeřáb} \textsuperscript{(32)}, který zvedá jako pírka obrovská břemena.

%%K t10-tit-5 sub
\VUsaut{3.0mm}
\VUpk{10.2pt}{40}{Lodivod.}
\VUsaut{3.0mm}

%%K t10-14-1 p
14. --- Obdivoval jsem \VUgras{zaoceánské parníky} \textsuperscript{(34)} vyplouvající z rejdy
se svými \VUgras{vlajkami} \textsuperscript{(35)} vlajícími ve větru, řízené obratným
\VUgras{lodivodem} \textsuperscript{(36)}, který uměl podivuhodně sledovat
\VUgras{plavební dráhu} vyznačenou \VUgras{bójemi} \textsuperscript{(40)} a
\VUgras{znaky}, vyhýbaje se \VUgras{nákladním plachetnicím} \textsuperscript{(41)},
které se namáhavě řídí svou jedinou čtvercovou plachtou. Rozeznával jsem na
\VUgras{přídi} \textsuperscript{(37)} --- na přední části --- \VUgras{kajuty}
\textsuperscript{(39)} druhé třídy a na \VUgras{zádi} \textsuperscript{(38)} --- na zadní
části --- salony pro cestující první třídy.

%%K t10-15-1 p
15. --- Někdy jsem byl také přítomen nakládání \VUgras{nákladní lodi} \textsuperscript{(42)},
v níž se právě hromadilo zboží ze \VUgras{skladiště} \textsuperscript{(43)} a z
\VUgras{námořního nádraží} \textsuperscript{(44)}, přivezené četnými vlaky
\VUgras{nábřežní tratě} \textsuperscript{(45)}.

%%K t10-tit-6 sub
\VUsaut{3.0mm}
\VUpk{10.2pt}{40}{Radost z veslování.}
\VUsaut{3.0mm}

%%K t10-16-1 p
16. --- Často jsem jezdil na loďce v \VUgras{plovoucím doku} \textsuperscript{(53)}, do něhož
se uchylují \VUgras{bárky} \textsuperscript{(46)}, a bylo mi radostí ovládat
\VUgras{veslo} \textsuperscript{(47)}. Vystoupil jsem na \VUgras{palubu}
\textsuperscript{(48)} \VUgras{plachetní bárky} \textsuperscript{(49)}, šplhal jsem po
\VUgras{lanech} \textsuperscript{(52)} na \VUgras{stěžeň} \textsuperscript{(51)} až do
\VUgras{ráhen} \textsuperscript{(50)}, potom jsem opakoval svou projížďku
\VUgras{na jole} \textsuperscript{(54)} mezi těžkými \VUgras{rybářskými bárkami}
\textsuperscript{(55)}, kde schnou \VUgras{sítě} \textsuperscript{(56)}.

%%K t10-17-1 p
17. --- Na \VUgras{nábřeží} \textsuperscript{(58)} přede mnou defilovalo nejrozmanitější
\VUgras{zboží} \textsuperscript{(59)}, uložené v \VUgras{bednách}
\textsuperscript{(60)} nebo v \VUgras{balících} \textsuperscript{(61)}. Tvořilo
\VUgras{náklad} \textsuperscript{(63)} \VUgras{nákladních člunů}, a mohutné
\VUgras{jeřáby} \textsuperscript{(62)} je zvedaly a skládaly na \VUgras{nákladní
vozy} \textsuperscript{(64)}, zatímco \VUgras{formani} podávali své prohlášení a platili
poplatky v \VUgras{celnici} \textsuperscript{(65)}.

%%K t10-18-1 p
18. --- Konečně, když jsem došel k \VUgras{otočnému mostu} \textsuperscript{(66)} nebo k
\VUgras{plavební komoře} \textsuperscript{(69)}, procházeje před \VUgras{bagrem}
\textsuperscript{(68)}, který čistí \VUgras{průplav} \textsuperscript{(67)}, vystoupil jsem na
břeh na molu blízko \VUgras{semaforu} \textsuperscript{(70)}.

%%K t10-19-1 p
19. --- Je to na druhé straně toho mola, kde přistávají \VUgras{šalupy}
\textsuperscript{(71)}, které přivážejí na břeh důstojníky eskadry. Je to obdivuhodná podívaná
vidět námořníky, jak v taktu zvedají svá vesla, zatímco bárka, nesená
\VUgras{vlnou} \textsuperscript{(72)}, snadno přistává u schůdků přístaviště. Na tomto
místě lze nejlépe pozorovat \VUgras{dmutí} a pohyb \VUgras{přílivu} a
\VUgras{odlivu} na \VUgras{pláži} \textsuperscript{(73)}.

%%K t10-tit-7 sub
\VUsaut{3.0mm}
\VUpk{10.2pt}{40}{Klub.}
\VUsaut{3.0mm}

%%K t10-20-1 p
20. --- K návratu domů jsem použil \VUgras{elektrickou tramvaj} \textsuperscript{(74)}, jejíž
\VUgras{zastávka} \textsuperscript{(75)} je blízko \VUgras{sloupu} postaveného před
důstojnickým klubem.

%%K t10-20-2 p
Z \VUgras{balkonu} \textsuperscript{(76)} klubu, ozdobeného \VUgras{kotvami}
\textsuperscript{(77)}, děly a koulemi, jsem často vídal skvělé shromáždění námořních
důstojníků: \VUgras{poručíky} \textsuperscript{(78)} s jejich kordem,
\VUgras{kapitány dělostřelectva} \textsuperscript{(79)} a \VUgras{pěchoty}
\textsuperscript{(80)} s jejich \VUgras{šavlí}. Za nimi stál \VUgras{námořník}
\textsuperscript{(81)}, který jim přinášel \VUgras{dalekohled}, když chtěli rozeznat, co se
děje v dálce.

%%K t10-21-1 p
21. --- Viděl jsem také mladé \VUgras{praporčíky} \textsuperscript{(82)}, jak přijímají
rozkazy od svého \VUgras{velitele} \textsuperscript{(83)} nebo od \VUgras{admirála}
\textsuperscript{(84)}, zatímco \VUgras{poddůstojník} \textsuperscript{(85)} čekal, aby je
odnesl na loď.

%%K t10-22-1 p
22. --- Z toho balkonu je nádherný pohled: ovládá se celá pláž s jejími \VUgras{stany}
\textsuperscript{(86)} a \VUgras{kabinami} \textsuperscript{(87)}, a je vidět, v hodinu
koupání, \VUgras{koupající se} \textsuperscript{(88)}, kteří vesele dovádějí před
\VUgras{kasinem} \textsuperscript{(89)}.

%%K t10-23-1 p
23. --- Na konci pláže tvoří pobřeží malý \VUgras{poloostrov} \textsuperscript{(91)}
\VUgras{skalnatý}, o který se tříští \VUgras{velká vlna} \textsuperscript{(92)} a na
němž je postaven \VUgras{maják} \textsuperscript{(93)}, který v noci ukazuje cestu plavcům.
Ten poloostrov je spojen s pevninou jen velmi úzkou \VUgras{šíjí} \textsuperscript{(90)}.

%%K t10-tit-8 sub
\VUsaut{3.0mm}
\VUpk{10.2pt}{40}{Celkový pohled na pobřeží.}
\VUsaut{3.0mm}

%%K t10-24-1 p
24. --- \VUgras{Útes} \textsuperscript{(94)} se táhne dále, rozerván mořem, které v něm
rozmarně kreslí řadu \VUgras{zátok} \textsuperscript{(95)} a \VUgras{výběžků} (mysů), a činí
je nanejvýš malebným.

%%K t10-24-2 p
Na \VUgras{náhorní plošině} \textsuperscript{(96)}, která ovládá \VUgras{úžinu}
\textsuperscript{(98)}, je velmi významná \VUgras{tvrz} \textsuperscript{(97)} a několik
«vil».

%%K t10-25-1 p
25. --- Ta úžina odděluje od pevniny velmi nebezpečný \VUgras{ostrov} \textsuperscript{(99)},
obklopený \VUgras{úskalími} \textsuperscript{(100)} a
\VUgras{příboji}, kde bárky a dokonce velké lodi někdy uváznou. Přes rychlost pomoci a
rychlý příjezd \VUgras{záchranných člunů} jsou ta \VUgras{ztroskotání} strašlivá a za
rovnodennostních \VUgras{bouří} zahynulo několik lodí s lidmi i se zbožím.

%%K t10-25-2 p
Přes to sousedství je \VUgras{přístav} velmi navštěvovaný námořníky.
